% 所有讲的笔记共用的导言区（版式设置）。每讲的主文件用 % 所有讲的笔记共用的导言区（版式设置）。每讲的主文件用 % 所有讲的笔记共用的导言区（版式设置）。每讲的主文件用 % 所有讲的笔记共用的导言区（版式设置）。每讲的主文件用 \input{preamble} 引入。
\PassOptionsToPackage{fontset=fandol}{ctex}
\documentclass[lang=cn,scheme=chinese,color=blue,mode=fancy,oneside,11pt]{elegantbook}

\setcounter{tocdepth}{1}

\lstdefinestyle{agentpython}{
  language=Python,
  basicstyle=\small\ttfamily,
  backgroundcolor=\color{structurecolor!5},
  keywordstyle=\color{structurecolor}\bfseries,
  commentstyle=\color{black!55}\itshape,
  stringstyle=\color{teal!75!black},
  showstringspaces=false,
  frame=single,
  rulecolor=\color{structurecolor!25},
  framesep=6pt,
  breaklines=true,
  columns=fullflexible
}

% 纯文本代码块（Markdown、提示词格式等）：不按 Python 高亮，避免 # 开头的行被当成注释
\lstdefinestyle{agenttext}{
  basicstyle=\small\ttfamily,
  backgroundcolor=\color{structurecolor!5},
  showstringspaces=false,
  frame=single,
  rulecolor=\color{structurecolor!25},
  framesep=6pt,
  breaklines=true,
  columns=fullflexible
}

% 段内小标题：加粗、不缩进、上方留一点空，正文另起一段
\newcommand{\keypoint}[1]{\par\medskip\noindent\textbf{#1}\par\nobreak\smallskip}
 引入。
\PassOptionsToPackage{fontset=fandol}{ctex}
\documentclass[lang=cn,scheme=chinese,color=blue,mode=fancy,oneside,11pt]{elegantbook}

\setcounter{tocdepth}{1}

\lstdefinestyle{agentpython}{
  language=Python,
  basicstyle=\small\ttfamily,
  backgroundcolor=\color{structurecolor!5},
  keywordstyle=\color{structurecolor}\bfseries,
  commentstyle=\color{black!55}\itshape,
  stringstyle=\color{teal!75!black},
  showstringspaces=false,
  frame=single,
  rulecolor=\color{structurecolor!25},
  framesep=6pt,
  breaklines=true,
  columns=fullflexible
}

% 纯文本代码块（Markdown、提示词格式等）：不按 Python 高亮，避免 # 开头的行被当成注释
\lstdefinestyle{agenttext}{
  basicstyle=\small\ttfamily,
  backgroundcolor=\color{structurecolor!5},
  showstringspaces=false,
  frame=single,
  rulecolor=\color{structurecolor!25},
  framesep=6pt,
  breaklines=true,
  columns=fullflexible
}

% 段内小标题：加粗、不缩进、上方留一点空，正文另起一段
\newcommand{\keypoint}[1]{\par\medskip\noindent\textbf{#1}\par\nobreak\smallskip}
 引入。
\PassOptionsToPackage{fontset=fandol}{ctex}
\documentclass[lang=cn,scheme=chinese,color=blue,mode=fancy,oneside,11pt]{elegantbook}

\setcounter{tocdepth}{1}

\lstdefinestyle{agentpython}{
  language=Python,
  basicstyle=\small\ttfamily,
  backgroundcolor=\color{structurecolor!5},
  keywordstyle=\color{structurecolor}\bfseries,
  commentstyle=\color{black!55}\itshape,
  stringstyle=\color{teal!75!black},
  showstringspaces=false,
  frame=single,
  rulecolor=\color{structurecolor!25},
  framesep=6pt,
  breaklines=true,
  columns=fullflexible
}

% 纯文本代码块（Markdown、提示词格式等）：不按 Python 高亮，避免 # 开头的行被当成注释
\lstdefinestyle{agenttext}{
  basicstyle=\small\ttfamily,
  backgroundcolor=\color{structurecolor!5},
  showstringspaces=false,
  frame=single,
  rulecolor=\color{structurecolor!25},
  framesep=6pt,
  breaklines=true,
  columns=fullflexible
}

% 段内小标题：加粗、不缩进、上方留一点空，正文另起一段
\newcommand{\keypoint}[1]{\par\medskip\noindent\textbf{#1}\par\nobreak\smallskip}
 引入。
\PassOptionsToPackage{fontset=fandol}{ctex}
\documentclass[lang=cn,scheme=chinese,color=blue,mode=fancy,oneside,11pt]{elegantbook}

\setcounter{tocdepth}{1}

\lstdefinestyle{agentpython}{
  language=Python,
  basicstyle=\small\ttfamily,
  backgroundcolor=\color{structurecolor!5},
  keywordstyle=\color{structurecolor}\bfseries,
  commentstyle=\color{black!55}\itshape,
  stringstyle=\color{teal!75!black},
  showstringspaces=false,
  frame=single,
  rulecolor=\color{structurecolor!25},
  framesep=6pt,
  breaklines=true,
  columns=fullflexible
}

% 纯文本代码块（Markdown、提示词格式等）：不按 Python 高亮，避免 # 开头的行被当成注释
\lstdefinestyle{agenttext}{
  basicstyle=\small\ttfamily,
  backgroundcolor=\color{structurecolor!5},
  showstringspaces=false,
  frame=single,
  rulecolor=\color{structurecolor!25},
  framesep=6pt,
  breaklines=true,
  columns=fullflexible
}

% 段内小标题：加粗、不缩进、上方留一点空，正文另起一段
\newcommand{\keypoint}[1]{\par\medskip\noindent\textbf{#1}\par\nobreak\smallskip}
